
\documentclass[tikz,border=2pt]{standalone}
\usepackage[utf8]{inputenc}
\usepackage[T1]{fontenc}
\usepackage[french]{babel}
\usepackage{amsmath,amssymb}
\usepackage[table]{xcolor}
\usepackage{siunitx}
\usepackage[version=4]{mhchem}
\usepackage{tabularx,booktabs,array}
\usepackage{tikz}
\usetikzlibrary{arrows.meta,positioning,calc,fit,decorations.pathmorphing,patterns,shapes.geometric}
\usepackage{pgfplots}
\pgfplotsset{compat=1.18}
\definecolor{primary}{RGB}{14,116,144}
\definecolor{accent}{RGB}{16,185,129}
\definecolor{dark}{RGB}{15,23,42}
\definecolor{bglight}{RGB}{241,245,249}
\definecolor{borderlight}{RGB}{203,213,225}
\definecolor{examplepink}{RGB}{236,72,153}
\definecolor{remarkblue}{RGB}{14,165,233}
\definecolor{notepurple}{RGB}{99,102,241}
\definecolor{synthgreen}{RGB}{5,150,105}
\definecolor{synthorange}{RGB}{249,115,22}
\definecolor{synthviolet}{RGB}{79,70,229}
\definecolor{primary}{RGB}{14,116,144}
\definecolor{accent}{RGB}{16,185,129}
\definecolor{dark}{RGB}{15,23,42}
\definecolor{bglight}{RGB}{241,245,249}
\definecolor{borderlight}{RGB}{203,213,225}
\definecolor{examplepink}{RGB}{236,72,153}
\definecolor{remarkblue}{RGB}{14,165,233}
\definecolor{notepurple}{RGB}{99,102,241}
\definecolor{synthgreen}{RGB}{5,150,105}
\definecolor{synthorange}{RGB}{249,115,22}
\definecolor{synthviolet}{RGB}{79,70,229}
\definecolor{primary}{RGB}{14,116,144}
\definecolor{accent}{RGB}{16,185,129}
\definecolor{dark}{RGB}{15,23,42}
\definecolor{bglight}{RGB}{241,245,249}
\definecolor{examplepink}{RGB}{236,72,153}
\definecolor{remarkblue}{RGB}{14,165,233}
\definecolor{notepurple}{RGB}{99,102,241}
\definecolor{synthgreen}{RGB}{5,150,105}
\definecolor{synthorange}{RGB}{249,115,22}
\definecolor{synthviolet}{RGB}{79,70,229}
\begin{document}
\begin{tikzpicture}[scale=0.6, >=Stealth]
% Axe optique
\draw[->] (-2,0) -- (15,0) node[right]{\small axe optique};
% Objectif
\draw[thick, blue] (0,-2.5) -- (0,2.5);
\node[above, blue] at (0,2.7) {$L_1$ (objectif)};
\node[below, blue] at (0,-2.7) {$O_1$};
% Foyers
\filldraw (5.5,0) circle (2pt);
\node[below] at (5.5,-0.3) {\small $F'_1 = F_2$};
% Oculaire
\draw[thick, red] (11,-1.5) -- (11,1.5);
\node[above, red] at (11,1.7) {$L_2$ (oculaire)};
\node[below, red] at (11,-1.7) {$O_2$};
% Foyer image oculaire
\filldraw[red] (13.5,0) circle (2pt);
\node[below, red] at (13.5,-0.3) {\small $F'_2$};
% Rayons entrants (quasi-parallèles, faible inclinaison)
\draw[->, thick, orange] (-2,1.0) -- (0,0.8);
\draw[->, thick, orange] (-2,0.3) -- (0,0.1);
% Convergence vers image intermédiaire
\draw[thick, orange] (0,0.8) -- (5.5,-0.5);
\draw[thick, orange] (0,0.1) -- (5.5,-0.5);
% Image intermédiaire
\draw[thick, green!60!black, ->] (5.5,0) -- (5.5,-0.5);
\node[right, green!60!black] at (5.6,-0.5) {\small $B_1$};
\node[above right, green!60!black] at (5.5,0.1) {\small $A_1$};
% Rayons vers oculaire
\draw[thick, purple] (5.5,-0.5) -- (11,-1.8);
\draw[thick, purple] (5.5,-0.5) -- (11,0);
% Rayons émergents
\draw[->, thick, purple] (11,-1.8) -- (14,-3.0);
\draw[->, thick, purple] (11,0) -- (14,-1.2);
% Angle theta
\node at (-1,0.7) {\small $\theta$};
% Angle theta'
\node at (13,-0.8) {\small $\theta'$};
\end{tikzpicture}
\end{document}
